\section{Individual amplitudes, inverse thresholds and cumulants}\label{sec:consequences}
\subsection{Proof of the full equivalent}
Write $R_n(s)=\Log[Z_n(s)/Z_n(c)]$ on $\HH_+$, normalized at $c$.
The exact adjacent identity and the now-proved companion limit give
\begin{equation}\label{adjacentsum}
 R_n(s)+R_{n+1}(s)=nf(s^2)+E_n(s^2)+\Log(s/c).
\end{equation}
The exact boundary definition gives
$R_{n+1}(s)-R_n(s)=\Log(s/c)-\Log F_n(s)$.
Because $F_n\to(s/c)e^{-f(s^2)/2}$ locally uniformly and all functions
are nonzero with the same normalization, their analytic logarithms
converge locally uniformly. Hence
\begin{equation}\label{adjacentdifference}
 R_{n+1}(s)-R_n(s)\longrightarrow\tfrac12f(s^2).
\end{equation}
Subtract \eqref{adjacentdifference} from \eqref{adjacentsum} and divide
by two. This yields
\begin{equation}\label{Rlimit}
 R_n(s)-\tfrac n2f(s^2)
 \longrightarrow\tfrac12E(s^2)+\tfrac12\Log(s/c)-\tfrac14f(s^2)
\end{equation}
locally uniformly. Adding the real full-sequence calibration
\eqref{calibrationequiv} proves Theorem~\ref{thm:main}, including its
complex local uniformity. All factors in \eqref{amplitude} are positive
on the positive axis and nonzero on $\HH_+$. There are no zeros or
poles of $P$ there; the definitions at $c$ fix all roots uniquely.
At $s=1$, $P(1)=2/3$, so
$(1/c)^{1/2}P(1)^{-1/4}=2^{-1/4}$. Also
$b+\tfrac12\log(2/3)=\beta$. This verifies \eqref{aequiv}.

\subsection{Amplitude-aware inverse count thresholds}
Let
\[
 C_*=C(1)=C_{\rm cal}2^{-1/4}e^{E(1)/2},\qquad
 N(T)=\min\{n\ge1:a_n\ge T\},\quad y=\log T,\quad r=\sqrt{y/\alpha}.
\]
The minimum exists because \eqref{aequiv} tends to infinity.
The counts are nondecreasing: $A\mapsto\diag(A,1)$ is an injection.
Only large $T$ is considered below.

\begin{corollary}\label{cor:inverse}
Define the real center
\begin{equation}\label{zinverse}
 z_C(T)=r-\frac{\beta}{2\alpha}
       +\frac{\kappa\log r+\beta^2/(4\alpha)-\log C_*}{2\alpha r}.
\end{equation}
There exists a nonnegative function $\eps(T)\to0$ such that
\begin{equation}\label{inversebracket}
 \ceil{z_C(T)-\eps(T)/r}\ \le N(T)\le\ \ceil{z_C(T)+\eps(T)/r}.
\end{equation}
Equivalently the displayed center is
\begin{align}\label{inverseexpanded}
 z_C(T)={}&\sqrt{\frac{\log T}{\alpha}}-\frac{\beta}{2\alpha}
 +\frac{\kappa\log\log T}{4\sqrt{\alpha\log T}}\nonumber\\
 &+\frac{\beta^2/(4\alpha)-\log C_*-(\kappa/2)\log\alpha}
             {2\sqrt{\alpha\log T}}.
\end{align}
Thus the amplitude enters at the $1/\sqrt{\log T}$ scale, and the
remaining error is little-o at that scale before ceilings.
\end{corollary}
\begin{proof}
Set $g(x)=\alpha x^2+\beta x-\kappa\log x+\log C_*$ for $x>0$.
The full equivalent says $\log a_n=g(n)+\eta_n$ with $\eta_n\to0$.
On $[r/2,2r]$, $g'(x)=2\alpha x+\beta-\kappa/x\asymp r$.
Substitution of \eqref{zinverse}, writing its last numerator as
$L_r=\kappa\log r+\beta^2/(4\alpha)-\log C_*$, gives
\[
 g(z_C(T))-y=O((\log r)/r)=o(1).
\]
Indeed the quadratic and linear contributions are
$y-\beta^2/(4\alpha)+L_r+O((\log r)^2/r^2)$,
while $-\kappa\log z_C=-\kappa\log r+O(1/r)$.
The leading terms cancel exactly.

The growth theorem first gives $N(T)\in[r/2,2r]$ for all sufficiently
large $T$ (the endpoints have different fixed quadratic multiples
of $y$). On this interval the tail supremum
$\sup_{n\ge r/2}|\eta_n|$ tends to zero.
Choose $\eps(T)>0$, tending to zero, larger by a fixed sufficient
factor than this supremum plus $|g(z_C)-y|$ and $1/r$.
The lower derivative bound then shows that every integer
$n<z_C-\eps/r$ has $\log a_n<y$, while the integer
$n=\ceil{z_C+\eps/r}$ has $\log a_n\ge y$.
Integers outside $[r/2,2r]$ are handled by the preliminary growth
comparison and monotonicity. These two statements are exactly
\eqref{inversebracket}. Substituting $r=\sqrt{y/\alpha}$ and
$\log r=(\log y-\log\alpha)/2$ proves \eqref{inverseexpanded}.
\end{proof}

The ceilings are essential: even a vanishing real error can decide
which adjacent integer is the threshold when the center is near an
integer. The theorem does not give a fixed explicit decay rate for
$\eps(T)$, nor does it justify discarding it inside a ceiling.

\subsection{Constant terms of every fixed diagonal cumulant}
For fixed $s>0$ assign probability $s^{S(A)}/Z_n(s)$ to a DSASM,
and let $S_n$ denote its number of nonzero diagonal entries.
Put $\mathcal D=s\,d/ds$ and $L(s)=\Log C(s)$, real on the
positive axis. Define
\[
 m(s)=\tfrac12-\frac1{(s+1)(s+2)},\qquad
 v(s)=\frac{s(2s+3)}{(s+1)^2(s+2)^2}>0.
\]
Then $\mathcal DB=m$ and $\mathcal Dm=v$.

\begin{corollary}\label{cor:cumulants}
For each fixed $k\ge1$, locally uniformly on positive fugacity intervals,
\begin{equation}\label{cumulantconstant}
 \operatorname{cum}_k(S_n)=n\mathcal D^{k-1}m(s)
                              +\mathcal D^kL(s)+o(1).
\end{equation}
In particular
\begin{align*}
 \mathbb E_sS_n&=nm(s)+s^2E'(s^2)+\tfrac12-\tfrac12m(s)+o(1),\\
 \Var_s(S_n)&=nv(s)+2s^2E'(s^2)+2s^4E''(s^2)-\tfrac12v(s)+o(1).
\end{align*}
At $s=1$ the linear coefficients are $m(1)=1/3$ and $v(1)=5/36$.
\end{corollary}
\begin{proof}
The locally holomorphic limit \eqref{Rlimit} implies, on a small
complex disk about $z=0$,
\[
 \log\mathbb E_se^{zS_n}
 =n[B(se^z)-B(s)]+L(se^z)-L(s)+o(1)
\]
locally uniformly. Cauchy's formula makes every fixed derivative
of the error $o(1)$. Differentiating proves
\eqref{cumulantconstant}. Inserting
$L(s)=\log C_{\rm cal}+E(s^2)/2+\tfrac12\Log(s/c)
-\tfrac14\Log P(s)$ gives the displayed first two constants.
\end{proof}
This strengthens the bounded cumulant errors in the earlier article.
It relies on complex-local convergence, rather than differentiating
pointwise real asymptotics.

\section{Scope and reproducibility}\label{sec:scope}
The result is a full leading equivalent at each fixed positive
fugacity, with a positive constant characterized by convergent
operators and a Fourier series. Convergence is locally uniform on
the right half-plane. It is not asserted uniformly as $s\to0$ or
$s\to\infty$, nor uniformly in an unbounded derivative order.
No numerical digits for $C(1)$, explicit rate for the little-o terms,
$n^{-2}$ correction, or all-orders conjecture are established here.

The source bundle contains this article's modular TeX, its PDF,
the unchanged earlier-input article's PDF and TeX, a closed SHA-256
inventory, deterministic clean-build and repack scripts, and exact
finite checks with an adversarial input campaign. The finite checks
corroborate identities and normalization; they are not proofs of any
infinite-dimensional estimate, asymptotic convergence or numerical
amplitude. The analytical arguments above establish those limits.
The replay checks package immutability, normal and optimized Python
behavior, independent exact data, two byte-identical clean PDF builds,
and a closed deterministic ZIP round trip. Rebuilding needs only the
locally installed TeX distribution and Python; the scripts make no
network requests and install no dependencies.
